\documentclass[../../main.tex]{subfiles}

\begin{document}

\subsection{Il magnete} \label{subsec:il-magnete}
  \begin{figure}[ht]
    \centering
    \includegraphics[width=.5\textwidth]{DUNE/sand/cal-mag}
    \caption{
      Vista frontale (figura in alto) e laterale (figura in basso)
      dei calorimetri, che mostra le guide di luce e la loro posizione
      all'interno del magnete di KLOE. Le unità di misura sono in mm.
      \cite{abud2021deepundergroundneutrinoexperiment}
    }
    \label{fig:cal-mag}
  \end{figure}
  Il magnete che verrà utilizzato in SAND è lo stesso dell'esperimento
  KLOE~\cite{Krzemien_2019}, esso produce un campo magnetico di 0,6~T in un
  volume cilindrico magnetizzato lungo 4,3~m e con un diametro di 4,8~m.

  Il raffreddamento avviene tramite elio: l'elio gassoso a 5,2~K e a una
  pressione di 3~bar viene liquefatto e posto in contatto termico con la bobina
  che opera a una corrente di 2902~A e l'energia immagazzinata è di 14.32~MJ\@.
  Lo schermo radiativo (\textit{radiation shield}) è invece raffreddato da elio
  gassoso a 70~K proveniente dall'impianto criogenico.
  Uno schema del magnete è riportato in figura~\ref{fig:cal-mag}.

\end{document}